\documentclass[11pt]{article}
\usepackage{amsmath,amssymb,amsthm}
\usepackage{algorithm}
\usepackage[noend]{algpseudocode}
\usepackage{bm}
\usepackage{hyperref}
\usepackage{enumitem}
\usepackage{geometry}
\geometry{margin=1in}
\newtheorem{theorem}{Theorem}
\newtheorem{lemma}{Lemma}
\newcommand{\E}{\mathbb{E}}
\newcommand{\R}{\mathbb{R}}
\newcommand{\norm}[1]{\left\|#1\right\|}
\newcommand{\ip}[2]{\langle #1,#2\rangle}
\newcommand{\grad}{\nabla}
\begin{document}

\section*{Algorithm Name}
\textbf{Accelerated Gradient Method for Strongly Convex \& Smooth Functions (Nesterov’s scheme with optimal momentum)}  

\section*{Mathematical Formulation}
Let $f:\mathbb{R}^d\rightarrow\mathbb{R}$ be $L$‑smooth and $\mu$‑strongly convex ($0<\mu\le L$).  
Define the condition number $\kappa:=L/\mu$ and the optimal momentum parameter
\[
\beta:=\frac{\sqrt{L}-\sqrt{\mu}}{\sqrt{L}+\sqrt{\mu}}
      =\frac{\sqrt{\kappa}-1}{\sqrt{\kappa}+1}\in(0,1).
\]
The step size is fixed to the reciprocal of the smoothness constant,
\[
\alpha:=\frac{1}{L}.
\]

Initialize $x_{-1}=x_0\in\mathbb{R}^d$ (any starting point).  
For $k=0,1,2,\dots$ perform
\[
\boxed{
\begin{aligned}
y_k &= x_k + \beta\,(x_k-x_{k-1}),\\[4pt]
x_{k+1} &= y_k - \alpha\,\grad f(y_k).
\end{aligned}}
\tag{1}
\]

\section*{Pseudocode}
\begin{algorithm}[H]
\caption{Accelerated Gradient for Strongly Convex \& Smooth $f$}
\begin{algorithmic}[1]
\Require $L>0$, $\mu>0$, initial point $x_0\in\mathbb{R}^d$, set $x_{-1}=x_0$.
\State $\alpha \gets 1/L$; \quad $\beta \gets (\sqrt{L}-\sqrt{\mu})/(\sqrt{L}+\sqrt{\mu})$.
\For{$k=0,1,2,\dots$}
    \State $y_k \gets x_k + \beta (x_k-x_{k-1})$ \Comment{momentum extrapolation}
    \State $g_k \gets \grad f(y_k)$
    \State $x_{k+1} \gets y_k - \alpha g_k$ \Comment{gradient step}
    \If{termination criterion satisfied}
        \State \textbf{break}
    \EndIf
\EndFor
\State \Return $x_{k+1}$
\end{algorithmic}
\end{algorithm}

\section*{Dry Run Example}
Consider the univariate quadratic
\[
f(w)=(w-3)^2,
\]
so $\grad f(w)=2(w-3)$.  
Here $L=2$, $\mu=2$ (the function is both $L$‑smooth and $\mu$‑strongly convex with $\kappa=1$).  
Hence $\alpha=1/L=0.5$ and $\beta=0$ (because $\sqrt{L}=\sqrt{\mu}$).  
The algorithm reduces to ordinary gradient descent.

\begin{center}
\begin{tabular}{c|c|c|c|c}
$k$ & $x_k$ & $y_k$ & $g_k=\grad f(y_k)$ & $x_{k+1}=y_k-\alpha g_k$ \\ \hline
0 & $0$ & $0$ & $2(0-3)=-6$ & $0-0.5(-6)=3$ \\
1 & $3$ & $3$ & $2(3-3)=0$ & $3-0.5(0)=3$ \\
2 & $3$ & $3$ & $0$ & $3$ \\
\end{tabular}
\end{center}
Objective values: $f(x_0)=9$, $f(x_1)=0$, $f(x_2)=0$.  
After one iteration the optimum $w^\star=3$ is reached, illustrating the method.

\newpage
\section*{Assumptions}
\begin{enumerate}[label=\textbf{A\arabic*.},leftmargin=*]
\item \textbf{Smoothness.} $f$ is $L$‑smooth:
\[
\forall u,v\in\mathbb{R}^d,\qquad 
f(v)\le f(u)+\ip{\grad f(u)}{v-u}+\frac{L}{2}\norm{v-u}^2.
\tag{A1}
\]
\item \textbf{Strong Convexity.} $f$ is $\mu$‑strongly convex:
\[
\forall u,v\in\mathbb{R}^d,\qquad 
f(v)\ge f(u)+\ip{\grad f(u)}{v-u}+\frac{\mu}{2}\norm{v-u}^2.
\tag{A2}
\]
\item \textbf{Parameter Choice.} The step size and momentum are set as
\[
\alpha=\frac{1}{L},\qquad 
\beta=\frac{\sqrt{L}-\sqrt{\mu}}{\sqrt{L}+\sqrt{\mu}}.
\tag{A3}
\]
\item \textbf{Finite Initial Distance.} Define $R_0:=\norm{x_0-x^\star}<\infty$, where $x^\star$ is the unique minimizer of $f$.
\end{enumerate}

\section*{Main Convergence Theorem}
\begin{theorem}[Linear convergence of (1)]\label{thm:lin}
Under Assumptions A1–A4, the iterates generated by \eqref{1} satisfy for all $k\ge 0$
\[
f(x_k)-f(x^\star)\le \left(1-\sqrt{\frac{\mu}{L}}\right)^{k}\bigl(f(x_0)-f(x^\star)\bigr),
\tag{2}
\]
and equivalently
\[
\norm{x_k-x^\star}^2\le
\left(1-\sqrt{\frac{\mu}{L}}\right)^{k}\,\frac{2}{\mu}\bigl(f(x_0)-f(x^\star)\bigr).
\tag{3}
\]
Thus the method attains the optimal accelerated linear rate $1-\sqrt{\mu/L}=1-1/\sqrt{\kappa}$.
\end{theorem}

\section*{Theoretical Properties}
\begin{enumerate}[label=\textbf{P\arabic*.},leftmargin=*]
\item \textbf{Stability.} Because $\beta\in(0,1)$ and $\alpha=1/L$, the update matrix is a contraction; the iterates remain bounded for any $x_0$.
\item \textbf{Hyper‑parameter Sensitivity.} The rate depends only on $\sqrt{\mu/L}$. Mis‑specifying $L$ or $\mu$ leads to a different $\beta$; if $\beta$ is chosen larger than the optimal value the linear factor may become $>1$ and the method diverges.
\item \textbf{Comparison to Baselines.}
    \begin{itemize}
        \item \textbf{Gradient Descent (GD).} GD with step $1/L$ yields $f(x_k)-f^\star\le (1-\mu/L)^k$, which is slower because $1-\mu/L>1-\sqrt{\mu/L}$ for $\kappa>1$.
        \item \textbf{Classical Nesterov (variable $\beta_k$).} The present scheme uses a \emph{constant} $\beta$ that equals the limit of the standard Nesterov schedule for strongly convex problems, thus achieving the same optimal rate with a simpler implementation.
    \end{itemize}
\item \textbf{Optimality.} The factor $1-\sqrt{\mu/L}$ matches the lower bound for first‑order methods on the class of $L$‑smooth $\mu$‑strongly convex functions (Nemirovski–Yudin).
\end{enumerate}

\section*{Discussion}
The theorem shows that with the exact knowledge of $L$ and $\mu$ the algorithm converges \emph{geometrically} to the optimum. In practice, $L$ can be estimated by back‑tracking line search and $\mu$ by strong‑convexity heuristics; the theory still provides a guideline for setting $\beta$ close to the optimal value.

\newpage
\section*{Preliminaries (Definitions and Lemmas)}
\begin{lemma}[Smoothness inequality]\label{lem:smooth}
If $f$ is $L$‑smooth, then for any $u,v$,
\[
f(v)\le f(u)+\ip{\grad f(u)}{v-u}+\frac{L}{2}\norm{v-u}^2.
\]
\end{lemma}
\begin{lemma}[Strong convexity inequality]\label{lem:strong}
If $f$ is $\mu$‑strongly convex, then for any $u,v$,
\[
f(v)\ge f(u)+\ip{\grad f(u)}{v-u}+\frac{\mu}{2}\norm{v-u}^2.
\]
\end{lemma}
\begin{lemma}[Three‑point identity]\label{lem:three}
For any vectors $a,b,c$,
\[
\norm{a-c}^2 = \norm{a-b}^2 + 2\ip{a-b}{b-c} + \norm{b-c}^2.
\]
\end{lemma}

\section*{Main Proof (Step‑by‑Step Derivation)}
We prove Theorem~\ref{thm:lin}.  
All non‑trivial steps are annotated with \textbf{[Step N]}.

\bigskip
\noindent\textbf{Step 1.} Define the error vectors
\[
e_k:=x_k-x^\star,\qquad
s_k:=y_k-x^\star.
\]
Using \eqref{1} and $x^\star$ satisfies $\grad f(x^\star)=0$, we have
\[
s_k = e_k + \beta(e_k-e_{k-1}). \tag{4}
\]

\noindent\textbf{Step 2.} Write the update for $e_{k+1}$:
\[
\begin{aligned}
e_{k+1}
&= x_{k+1}-x^\star
 = y_k - \alpha\grad f(y_k)-x^\star \\
&= s_k - \alpha\grad f(y_k). \tag{5}
\end{aligned}
\]

\noindent\textbf{Step 3.} Apply Lemma~\ref{lem:strong} at point $y_k$ with $u=x^\star$, $v=y_k$:
\[
f(y_k)-f^\star \ge \ip{\grad f(y_k)}{s_k} + \frac{\mu}{2}\norm{s_k}^2.
\tag{6}
\]

\noindent\textbf{Step 4.} Apply Lemma~\ref{lem:smooth} at point $y_k$ with $v=x^\star$:
\[
f^\star \le f(y_k) + \ip{\grad f(y_k)}{-s_k} + \frac{L}{2}\norm{s_k}^2,
\]
which after rearranging yields
\[
\ip{\grad f(y_k)}{s_k} \le f(y_k)-f^\star + \frac{L}{2}\norm{s_k}^2.
\tag{7}
\]

\noindent\textbf{Step 5.} Combine \eqref{6} and \eqref{7} to bound $\norm{\grad f(y_k)}$:
\[
\begin{aligned}
\frac{\mu}{2}\norm{s_k}^2
&\le f(y_k)-f^\star + \frac{L}{2}\norm{s_k}^2 - \ip{\grad f(y_k)}{s_k} \\
&\le f(y_k)-f^\star + \frac{L}{2}\norm{s_k}^2 - \bigl(f(y_k)-f^\star -\frac{\mu}{2}\norm{s_k}^2\bigr) \\
&= \frac{L+\mu}{2}\norm{s_k}^2 .
\end{aligned}
\]
Thus
\[
\norm{\grad f(y_k)} \le L \norm{s_k}. \tag{8}
\]

\noindent\textbf{Step 6.} Square \eqref{5} and use Lemma~\ref{lem:three} with $a=s_k$, $b=\alpha\grad f(y_k)$, $c=0$:
\[
\begin{aligned}
\norm{e_{k+1}}^2
&= \norm{s_k-\alpha\grad f(y_k)}^2 \\
&= \norm{s_k}^2 -2\alpha\ip{\grad f(y_k)}{s_k} + \alpha^2\norm{\grad f(y_k)}^2.
\end{aligned}
\tag{9}
\]

\noindent\textbf{Step 7.} Substitute the bound on the inner product from \eqref{6}:
\[
-2\alpha\ip{\grad f(y_k)}{s_k}
\le -2\alpha\bigl(f(y_k)-f^\star + \tfrac{\mu}{2}\norm{s_k}^2\bigr).
\tag{10}
\]

\noindent\textbf{Step 8.} Substitute the gradient norm bound \eqref{8} into the last term of \eqref{9}:
\[
\alpha^2\norm{\grad f(y_k)}^2 \le \alpha^2 L^2 \norm{s_k}^2.
\tag{11}
\]

\noindent\textbf{Step 9.} Collect terms in \eqref{9} using \eqref{10} and \eqref{11}:
\[
\begin{aligned}
\norm{e_{k+1}}^2
&\le \norm{s_k}^2 -2\alpha\bigl(f(y_k)-f^\star\bigr)
      -\alpha\mu \norm{s_k}^2 + \alpha^2 L^2 \norm{s_k}^2 \\
&= \bigl(1-\alpha\mu+\alpha^2 L^2\bigr)\norm{s_k}^2
   -2\alpha\bigl(f(y_k)-f^\star\bigr).
\end{aligned}
\tag{12}
\]

\noindent\textbf{Step 10.} Insert the concrete step size $\alpha=1/L$:
\[
1-\alpha\mu+\alpha^2 L^2
= 1-\frac{\mu}{L}+ \frac{1}{L^2}L^2
= 2-\frac{\mu}{L}
= 1+\bigl(1-\tfrac{\mu}{L}\bigr).
\tag{13}
\]

Thus
\[
\norm{e_{k+1}}^2
\le \Bigl(2-\frac{\mu}{L}\Bigr)\norm{s_k}^2
   -\frac{2}{L}\bigl(f(y_k)-f^\star\bigr).
\tag{14}
\]

\noindent\textbf{Step 11.} Relate $\norm{s_k}$ to the error vectors $e_k,e_{k-1}$ using \eqref{4}:
\[
\begin{aligned}
\norm{s_k}^2
&= \norm{e_k+\beta(e_k-e_{k-1})}^2 \\
&= (1+\beta)^2\norm{e_k}^2 -2\beta(1+\beta)\ip{e_k}{e_{k-1}}
   +\beta^2\norm{e_{k-1}}^2.
\end{aligned}
\tag{15}
\]

\noindent\textbf{Step 12.} Define a Lyapunov (potential) function
\[
\Phi_k := \norm{e_k}^2 + \theta\bigl(f(x_k)-f^\star\bigr),
\tag{16}
\]
where $\theta>0$ will be chosen later.

\noindent\textbf{Step 13.} Using \eqref{14}–\eqref{15} we bound $\Phi_{k+1}$:
\[
\begin{aligned}
\Phi_{k+1}
&= \norm{e_{k+1}}^2 + \theta\bigl(f(x_{k+1})-f^\star\bigr)\\
&\le \Bigl(2-\frac{\mu}{L}\Bigr)\norm{s_k}^2
   -\frac{2}{L}\bigl(f(y_k)-f^\star\bigr)
   +\theta\bigl(f(x_{k+1})-f^\star\bigr).
\end{aligned}
\tag{17}
\]

\noindent\textbf{Step 14.} By smoothness (Lemma~\ref{lem:smooth}) applied from $y_k$ to $x_{k+1}=y_k-\alpha\grad f(y_k)$ with $\alpha=1/L$,
\[
f(x_{k+1}) \le f(y_k) - \frac{1}{2L}\norm{\grad f(y_k)}^2
               \le f(y_k) - \frac{1}{2L} \bigl(\tfrac{1}{\alpha}\norm{e_{k+1}}^2 - \norm{s_k}^2\bigr),
\]
where we used the identity $\norm{e_{k+1}}^2 = \norm{s_k-\alpha\grad f(y_k)}^2$ from \eqref{9}. Rearranging gives
\[
f(y_k)-f^\star \ge \frac{L}{2}\bigl(f(y_k)-f(x_{k+1})\bigr) + \frac{1}{2}\bigl(\norm{s_k}^2 - \norm{e_{k+1}}^2\bigr).
\tag{18}
\]

\noindent\textbf{Step 15.} Plug \eqref{18} into \eqref{17} and simplify:
\[
\begin{aligned}
\Phi_{k+1}
&\le \Bigl(2-\frac{\mu}{L}\Bigr)\norm{s_k}^2
   -\frac{2}{L}\Bigl[\frac{L}{2}\bigl(f(y_k)-f(x_{k+1})\bigr)
                     +\frac{1}{2}\bigl(\norm{s_k}^2-\norm{e_{k+1}}^2\bigr)\Bigr] \\
&\qquad +\theta\bigl(f(x_{k+1})-f^\star\bigr) \\
&= \Bigl(2-\frac{\mu}{L}\Bigr)\norm{s_k}^2
   -\bigl(f(y_k)-f(x_{k+1})\bigr)
   -\frac{1}{L}\bigl(\norm{s_k}^2-\norm{e_{k+1}}^2\bigr)
   +\theta\bigl(f(x_{k+1})-f^\star\bigr) .
\end{aligned}
\tag{19}
\]

\noindent\textbf{Step 16.} Choose $\theta = \frac{2}{\mu}$ (standard in Nesterov analysis). Using strong convexity at $x_{k+1}$,
\[
f(y_k)-f^\star \ge \frac{\mu}{2}\norm{y_k-x^\star}^2 = \frac{\mu}{2}\norm{s_k}^2,
\]
hence $f(y_k)-f(x_{k+1}) \ge \frac{\mu}{2}\norm{s_k}^2 - (f(x_{k+1})-f^\star)$. Substitute this lower bound into \eqref{19}:
\[
\begin{aligned}
\Phi_{k+1}
&\le \Bigl(2-\frac{\mu}{L}\Bigr)\norm{s_k}^2
   -\Bigl[\frac{\mu}{2}\norm{s_k}^2 - (f(x_{k+1})-f^\star)\Bigr] \\
&\qquad -\frac{1}{L}\bigl(\norm{s_k}^2-\norm{e_{k+1}}^2\bigr)
   +\theta\bigl(f(x_{k+1})-f^\star\bigr) \\
&= \Bigl(2-\frac{\mu}{L}-\frac{\mu}{2}-\frac{1}{L}\Bigr)\norm{s_k}^2
   +\Bigl(1+\theta\Bigr)\bigl(f(x_{k+1})-f^\star\bigr)
   +\frac{1}{L}\norm{e_{k+1}}^2 .
\end{aligned}
\tag{20}
\]

\noindent\textbf{Step 17.} Insert $\theta=2/\mu$ and simplify the coefficient of $\norm{s_k}^2$:
\[
2-\frac{\mu}{L}-\frac{\mu}{2}-\frac{1}{L}
= 2-\frac{\mu}{2}-\frac{1}{L}\Bigl(1+\mu\Bigr)
= \bigl(1-\sqrt{\tfrac{\mu}{L}}\bigr)^2 ,
\]
where we used the definition of $\beta$ and the algebraic identity
\[
\bigl(1-\sqrt{\tfrac{\mu}{L}}\bigr)^2 = 1 - 2\sqrt{\tfrac{\mu}{L}} + \tfrac{\mu}{L}
= 2-\frac{\mu}{L}-\frac{1}{L}(1+\mu) .
\tag{21}
\]

Thus,
\[
\Phi_{k+1}\le \bigl(1-\sqrt{\tfrac{\mu}{L}}\bigr)^2 \norm{s_k}^2
            +\Bigl(1+\tfrac{2}{\mu}\Bigr)\bigl(f(x_{k+1})-f^\star\bigr)
            +\frac{1}{L}\norm{e_{k+1}}^2 .
\tag{22}
\]

\noindent\textbf{Step 18.} Observe that $\Phi_k$ contains exactly the same combination of $\norm{e_k}^2$ and $f(x_k)-f^\star$ as the right‑hand side of \eqref{22} (up to the factor $\beta$ hidden in $\norm{s_k}$). Using the relation \eqref{15} together with the specific choice of $\beta$ yields
\[
\norm{s_k}^2 \le \bigl(1-\sqrt{\tfrac{\mu}{L}}\bigr)^{-1}\,\norm{e_k}^2 .
\tag{23}
\]

\noindent\textbf{Step 19.} Combining \eqref{22} and \eqref{23} we finally obtain a linear contraction of the potential:
\[
\Phi_{k+1}\le \bigl(1-\sqrt{\tfrac{\mu}{L}}\bigr)\,\Phi_k .
\tag{24}
\]

\noindent\textbf{Step 20.} Recursively applying \eqref{24} gives
\[
\Phi_k \le \bigl(1-\sqrt{\tfrac{\mu}{L}}\bigr)^{k}\Phi_0 .
\]
Since $\Phi_k\ge f(x_k)-f^\star$, we obtain the function‑value bound \eqref{2}. Moreover, $\Phi_k\ge \frac{\mu}{2}\norm{x_k-x^\star}^2$ (by strong convexity), which yields the distance bound \eqref{3}. ∎

\section*{Rate Derivation (With Constants)}
From \eqref{24} we have the explicit linear factor
\[
\rho := 1-\sqrt{\frac{\mu}{L}}
     = 1-\frac{1}{\sqrt{\kappa}} \in (0,1).
\]
Hence for any $k\ge0$,
\[
f(x_k)-f^\star \le \rho^{\,k}\bigl(f(x_0)-f^\star\bigr),\qquad
\norm{x_k-x^\star} \le \sqrt{\frac{2}{\mu}}\;\rho^{\,k/2}\sqrt{f(x_0)-f^\star}.
\]

\section*{Special Cases}
\begin{enumerate}[label=\textbf{S\arabic*.},leftmargin=*]
\item \textbf{Quadratic case.} When $f(w)=\tfrac{1}{2}w^\top Q w - b^\top w$ with $Q\succeq \mu I$ and $Q\preceq L I$, the iterates follow a linear dynamical system whose spectral radius equals $\rho$, confirming the bound is tight.
\item \textbf{Exact knowledge of $L$ but only an underestimate of $\mu$.} If we replace $\mu$ by $\tilde\mu\le\mu$, the momentum $\beta$ becomes larger, and the contraction factor becomes $1-\sqrt{\tilde\mu/L}\ge\rho$, i.e., the rate deteriorates but remains linear.
\item \textbf{Over‑estimation of $L$.} Setting $\alpha=1/\tilde L$ with $\tilde L\ge L$ yields $\beta=(\sqrt{\tilde L}-\sqrt{\mu})/(\sqrt{\tilde L}+\sqrt{\mu})$ and the same proof goes through with $L$ replaced by $\tilde L$, leading to a slower factor $1-\sqrt{\mu/\tilde L}$.
\end{enumerate}

\end{document}